\clearpage
\section[\MakeUppercase{Эксперименты, метрики и анализ результатов}]{\centering\MakeUppercase{Эксперименты, метрики и анализ результатов}}

\subsection{\centering{Методика эксперимента}}

Экспериментальная часть основана на уже сохраненных результатах из каталога \code{results}. Для каждого решателя и датасета в JSON-файлах содержатся ответы по отдельным задачам, рассчитанные метрики и время решения. Такой подход выбран намеренно: он исключает повторную случайную генерацию результатов и позволяет опираться на фактические данные проекта.

При сравнении эвристических методов особенно важно использовать единый протокол оценки. Если один алгоритм проверяется более мягким валидатором, а другой -- более строгим, итоговое сравнение теряет смысл. В данном проекте все результаты проходят через \code{validator.py}, поэтому различия в score отражают именно свойства построенных укладок. Это также упрощает воспроизводимость: любой сохраненный JSON-ответ можно повторно проверить тем же валидатором и получить те же компоненты метрики.

Сохраненные результаты рассматриваются как экспериментальная выборка, а не как единичные демонстрационные примеры. Каждый датасет содержит множество задач с разным составом товаров, поэтому средние значения позволяют оценивать устойчивость метода к изменению входных данных. При этом отдельные наборы отличаются сложностью: \code{dataset\_benchmark.json} содержит больше объектов в задаче и более разнообразный сценарный состав, поэтому он используется как наиболее показательный контрольный набор.

В основном сравнении используются четыре single-pallet метода:
\begin{itemize}
  \item \code{greedy\_algorithm};
  \item \code{conditions} -- MaxRects-подобный послойный решатель;
  \item \code{gan\_ga\_solver};
  \item \code{lns\_solver}.
\end{itemize}

Общими для всех четырех решателей являются три набора: \code{dataset\_100.json}, \code{dataset\_200\_2.json} и \code{dataset\_benchmark.json}. Для GAN+GA дополнительно есть результат на \code{dataset\_150.json}, но он не используется в основном ранжировании, чтобы сравнение было симметричным.

Симметричность сравнения важна из-за того, что итоговый score усредняется по задачам внутри набора. Если включить датасет, на котором присутствует только один решатель, средний результат нельзя будет напрямую сопоставлять с остальными методами. Поэтому дополнительный результат GAN+GA полезен для анализа поведения метода, но не должен менять основную таблицу ранжирования.

\subsection{\centering{Сводные результаты}}

Средний итоговый score на общих наборах приведен в таблице~\ref{tab:score-main} и на рис.~\ref{fig:all-score}.

\begin{table}[H]
\caption{Средний итоговый score на общих наборах}
\label{tab:score-main}
\centering
\small
\begin{tabular}{|l|c|c|c|}
\hline
\textbf{Метод} & \textbf{\code{dataset\_100}} & \textbf{\code{dataset\_200\_2}} & \textbf{\code{dataset\_benchmark}} \\
\hline
Greedy & 0.6842 & 0.7120 & 0.6549 \\ \hline
MaxRects & 0.7104 & 0.7117 & 0.7124 \\ \hline
GAN+GA & 0.7012 & 0.7327 & 0.7161 \\ \hline
LNS & \textbf{0.7111} & \textbf{0.7434} & \textbf{0.7298} \\ \hline
\end{tabular}
\end{table}

\begin{figure}[H]
\centering
\includegraphics[width=0.97\textwidth]{all_solvers_score.png}
\caption{Средний итоговый score по решателям и датасетам}
\label{fig:all-score}
\end{figure}

LNS показывает лучший результат на всех трех общих наборах. На \code{dataset\_100.json} его преимущество над MaxRects минимально, что объясняется меньшей сложностью задач: быстрый послойный метод уже строит достаточно сильные решения. На \code{dataset\_200\_2.json} LNS заметно обгоняет остальные методы за счет лучшей полноты размещения. На \code{dataset\_benchmark.json}, где среднее число объектов выше и состав сценариев шире, преимущество LNS сохраняется.

Интерпретация этой таблицы показывает, что выигрыш LNS не является случайным эффектом одного датасета. Метод не всегда максимизирует отдельную компоненту, но стабильно удерживает высокий итоговый баланс. Для практической паллетизации такое поведение предпочтительнее, чем резкий выигрыш по одной характеристике при сильном проигрыше по другой. Например, алгоритм с максимальной утилизацией объема может оказаться менее полезным, если он часто размещает тяжелые товары над хрупкими или оставляет большую долю заказа вне паллеты.

Также важно учитывать, что различие между методами зависит от сложности набора. На простых задачах сильный baseline может быть близок к LNS, потому что пространство хороших решений достаточно широкое. На более насыщенных задачах ранние ошибки greedy и послойных методов становятся дороже: один неудачно выбранный нижний слой ограничивает все последующие уровни. В этих случаях возможность разрушить и восстановить часть укладки дает LNS преимущество.

\subsection{\centering{Компоненты метрики}}

Для интерпретации результатов рассмотрим компоненты score на контрольном наборе \code{dataset\_benchmark.json}. Данные приведены в таблице~\ref{tab:components-benchmark} и на рис.~\ref{fig:components}.

\begin{table}[H]
\caption{Компоненты метрики на \code{dataset\_benchmark.json}}
\label{tab:components-benchmark}
\centering
\small
\begin{tabular}{|l|c|c|c|c|c|}
\hline
\textbf{Метод} & \textbf{Score} & \textbf{Vol.} & \textbf{Coverage} & \textbf{Fragility} & \textbf{Time} \\
\hline
Greedy & 0.6549 & 0.6155 & 0.4906 & 1.0000 & 1.0000 \\ \hline
MaxRects & 0.7124 & \textbf{0.7513} & \textbf{0.6160} & 0.5198 & 1.0000 \\ \hline
GAN+GA & 0.7161 & 0.7156 & 0.5540 & \textbf{1.0000} & 0.9205 \\ \hline
LNS & \textbf{0.7298} & 0.7172 & 0.6036 & 0.9833 & 0.9175 \\ \hline
\end{tabular}
\end{table}

\begin{figure}[H]
\centering
\includegraphics[width=0.98\textwidth]{benchmark_components.png}
\caption{Компоненты score на контрольном наборе}
\label{fig:components}
\end{figure}

Greedy сохраняет идеальные значения по хрупкости и времени, но значительно уступает по покрытию заказа. MaxRects достигает максимальной утилизации объема и максимального покрытия на контрольном наборе, однако теряет почти половину балла за хрупкость. GAN+GA хорошо соблюдает хрупкость и заметно улучшает объем по сравнению с greedy, но уступает LNS по покрытию.

LNS не максимизирует отдельную компоненту любой ценой. Его сила заключается в балансе: утилизация объема близка к GAN+GA, покрытие близко к MaxRects, хрупкость остается высокой, а временной балл остается приемлемым. Именно этот баланс дает лучший итоговый score.

Компонентный анализ показывает, почему итоговая метрика полезнее одного показателя заполнения объема. MaxRects достигает наибольшей утилизации и покрытия на контрольном наборе, но сильная потеря по хрупкости снижает итоговую оценку. Greedy, наоборот, сохраняет идеальную хрупкость и скорость, но слишком часто оставляет часть заказа неразмещенной. LNS занимает промежуточное положение по отдельным компонентам, но именно поэтому оказывается наиболее устойчивым по суммарному критерию.

Такой результат согласуется с логикой выбранной метрики. В реальной задаче паллетизации нельзя оптимизировать только геометрию, игнорируя свойства товара. Укладка должна быть плотной, но не должна создавать чрезмерный риск повреждения, а время расчета должно оставаться приемлемым для использования в операционном контуре. Поэтому метод, который немного уступает по объему, но лучше соблюдает остальные требования, может получать более высокий итоговый score.

\subsection{\centering{Время выполнения}}

Среднее фактическое время решения приведено в таблице~\ref{tab:runtime}. Все методы на основных наборах остаются в пределах практически приемлемого времени. Greedy и MaxRects работают быстрее всех, LNS и GAN+GA тратят около секунды на контрольном наборе.

\begin{table}[H]
\caption{Среднее время решения, мс}
\label{tab:runtime}
\centering
\small
\begin{tabular}{|l|c|c|c|}
\hline
\textbf{Метод} & \textbf{\code{dataset\_100}} & \textbf{\code{dataset\_200\_2}} & \textbf{\code{dataset\_benchmark}} \\
\hline
Greedy & 114.8 & 118.5 & 139.0 \\ \hline
MaxRects & 15.1 & 15.6 & 52.0 \\ \hline
GAN+GA & 795.3 & 770.4 & 1080.6 \\ \hline
LNS & 474.3 & 523.6 & 1062.6 \\ \hline
\end{tabular}
\end{table}

Увеличение времени LNS оправдано ростом основных компонент качества. На контрольном наборе LNS теряет часть временного балла, но выигрывает по итоговой оценке. Это соответствует назначению метаэвристики: использовать дополнительное время не на случайное усложнение, а на улучшение структуры укладки.

С точки зрения прикладного применения абсолютное время LNS остается небольшим: на контрольном наборе оно измеряется примерно одной секундой. Для пакетного планирования паллет или предварительного расчета заказов такой уровень затрат обычно приемлем. Если же система должна работать строго в интерактивном режиме, можно уменьшить бюджет LNS и использовать greedy или MaxRects как быстрый fallback, сохраняя единый формат ответа и валидации.

При этом время выполнения следует рассматривать вместе с размером задачи и требуемым качеством ответа. Для небольших заказов выигрыш сложной метаэвристики может быть несущественным, поскольку быстрый решатель уже находит близкое по качеству решение. Для насыщенных заказов с большим числом коробов дополнительная секунда расчета может окупаться за счет более полного размещения и меньшего числа переносов товара на следующую паллету. Поэтому практическая система может выбирать метод адаптивно: запускать быстрый baseline для простых случаев и включать LNS для задач, где требуется улучшить плотность или покрытие.

\subsection{\centering{Сравнение LNS и Greedy}}

В рамках НИР отдельно исследовался LNS как развитие базового greedy. Для этой пары доступны графики, построенные из сохраненных результатов. На рис.~\ref{fig:lns-greedy-score} показано, что LNS стабильно превосходит greedy по среднему score на трех наборах.

\begin{figure}[H]
\centering
\includegraphics[width=0.85\textwidth]{lns_vs_greedy_score.png}
\caption{Сравнение среднего score для LNS и greedy}
\label{fig:lns-greedy-score}
\end{figure}

Компонентное сравнение на рис.~\ref{fig:lns-greedy-metrics} показывает, что главный вклад LNS дает по утилизации объема и покрытию заказа. Небольшое ухудшение по времени и хрупкости не перекрывает этот выигрыш.

\begin{figure}[H]
\centering
\includegraphics[width=0.95\textwidth]{lns_vs_greedy_metrics.png}
\caption{Сравнение компонент метрики для LNS и greedy}
\label{fig:lns-greedy-metrics}
\end{figure}

\subsection{\centering{Сравнение с brute force на малых задачах}}

Каталог \code{brute\_force\_vs\_algo} содержит отдельные эксперименты для малых и средних наборов. Их цель -- оценить поведение LNS рядом с более полным поиском. Результаты приведены в таблице~\ref{tab:brute}.

\begin{table}[H]
\caption{Сравнение brute force и LNS на специальных наборах}
\label{tab:brute}
\centering
\small
\begin{tabular}{|l|l|r|r|r|}
\hline
\textbf{Набор} & \textbf{Метод} & \textbf{Score} & \textbf{Coverage} & \textbf{Время, мс} \\
\hline
small & Brute force & 0.5563 & 0.9897 & 19.5 \\ \hline
small & LNS & \textbf{0.5602} & \textbf{1.0000} & 28.3 \\ \hline
medium & Brute force & 0.5681 & 0.9921 & 10440.7 \\ \hline
medium & LNS & \textbf{0.5780} & \textbf{1.0000} & 72.5 \\ \hline
hard & Brute force / LDS & 0.5886 & 0.9364 & 8815.0 \\ \hline
hard & LNS & \textbf{0.6501} & \textbf{1.0000} & 243.2 \\ \hline
\end{tabular}
\end{table}

На этих наборах LNS не только сопоставим с перебором, но и превосходит сохраненные результаты brute force/LDS по итоговой метрике. Важно, что для medium и hard время перебора резко возрастает, тогда как LNS остается в пределах сотен миллисекунд. Это подтверждает практическую ценность метаэвристического подхода.

\subsection{\centering{Многопаллетное расширение}}

В каталоге \code{multipallet} реализовано расширение на несколько паллет. Эта постановка ближе к реальной логистике: если весь заказ нельзя разместить на одной паллете, требуется распределить оставшиеся товары на новые паллеты. В проекте есть single-pallet и multi-pallet решения для greedy и LNS.

В рамках данной ВКР многопаллетный режим рассматривается как направление развития, поскольку основная метрика и основные benchmark-результаты сосредоточены на single-pallet постановке. Тем не менее наличие отдельного каталога \code{multipallet} показывает, что архитектура проекта допускает переход от локальной задачи плотной укладки к задаче распределения заказа между несколькими контейнерами.

Многопаллетная постановка меняет смысл оптимизации. В single-pallet варианте часть товара может остаться неразмещенной, и это отражается в компоненте покрытия. В multi-pallet режиме требуется решить, сколько паллет использовать и как распределить товары между ними. Тогда к текущим критериям добавляется стоимость дополнительной паллеты, а также могут появиться требования группировки товаров по магазинам, маршрутам или температурным зонам. Поэтому имеющееся расширение следует рассматривать как основу, которую нужно дополнить отдельной метрикой и отдельными экспериментами.

\subsection{\centering{Анализ ошибок и ограничений}}

Анализ кода и результатов позволяет выделить несколько ограничений текущей реализации.

Во-первых, модель устойчивости приближенная. Критерий 60\% опоры учитывает площадь контакта, но не анализирует центр масс, распределение давления, трение и деформации упаковки. Поэтому решение может быть валидным по формальному критерию, но требовать дополнительной инженерной проверки перед промышленным применением.

Во-вторых, синтетические данные не покрывают всего разнообразия реальных заказов. В генераторе есть полезные сценарии фуд-ритейла, но нет, например, температурных зон, запрета соседства, приоритетов выгрузки, ограничений по слоям и требований к группировке товаров по магазинам.

В-третьих, GAN+GA является облегченной исследовательской реализацией. Он использует NumPy-модель вместо полноценного обучения на большом корпусе решений. Поэтому текущие результаты не следует трактовать как окончательный предел ML-подходов к задаче.

В-четвертых, benchmark хранит средние значения, но не содержит статистического анализа устойчивости стохастических методов по нескольким независимым запускам на каждом датасете. Для LNS и GAN+GA в дальнейшем полезно оценивать дисперсию результата.

Еще одно ограничение связано с синтетической природой основной части данных. Генератор позволяет управлять сценариями и получать воспроизводимые наборы, однако реальные складские заказы могут иметь дополнительные зависимости: совместимость товарных групп, требования к последовательности выгрузки, ограничения по температуре, запрет соседства некоторых позиций и нестандартные размеры упаковок. Поэтому полученные результаты показывают качество методов в моделируемой постановке, но для промышленного внедрения их следует дополнительно проверять на реальных данных конкретной логистической системы.

Кроме того, сохраненные benchmark-результаты фиксируют итоговые средние значения, но не всегда показывают распределение ошибок по отдельным задачам. Для дальнейшего анализа полезно выделять группы экземпляров, на которых метод чаще всего проигрывает: заказы с большим числом мелких коробов, смеси тяжелых и хрупких товаров, задачи с ограниченными поворотами или случаи почти полного заполнения паллеты. Такая детализация поможет не только улучшить средний score, но и понять, какие именно ограничения требуют новых эвристик.

\subsection{\centering{Направления развития}}

Наиболее перспективные направления дальнейшей работы:
\begin{enumerate}
  \item Улучшить физическую модель устойчивости с учетом центра масс и распределения веса;
  \item Расширить генератор данных реальными бизнес-ограничениями;
  \item Добавить статистический протокол повторных запусков стохастических методов;
  \item Развить LNS за счет адаптивного выбора destroy/repair-операторов;
  \item Исследовать reinforcement learning или imitation learning для выбора порядка размещения;
  \item Развить многопаллетную постановку с метрикой минимизации числа паллет;
  \item Добавить автоматическую генерацию HTML-отчета с интерактивной 3D-визуализацией.
\end{enumerate}

\subsection{\centering{Выводы по разделу}}

Экспериментальные результаты показывают, что лучшим single-pallet методом текущего проекта является LNS. Он стабильно превосходит greedy, MaxRects и GAN+GA по среднему итоговому score на общих наборах. Причина преимущества не в максимизации одной компоненты, а в сбалансированном улучшении объема, покрытия, хрупкости и времени. Greedy остается полезным быстрым baseline, MaxRects -- сильным послойным решателем по плотности, GAN+GA -- перспективным исследовательским направлением, а brute force -- инструментом проверки малых экземпляров.
